%%============================================================================
%% thesis.tex  --  Adelaide Punt, DPhil thesis (body)
%%
%% Title: Noise and Uncertainty in Molecular Property Prediction:
%%        Robust Models and Adaptive Representations
%%============================================================================


\chapter{Introduction}\label{chap:intro}
% ---- TAKEAWAYS ----
% - A property prediction model is a molecular representation paired with a model; both limit what it can learn.
% - Labels in molecular data are noisy, and that noise is rarely homoscedastic Gaussian: its size varies with source and assay.
% - Under label noise, model architecture matters more than representation for robustness, and predicted uncertainty rises with noise but does not find the corrupted labels.
% - Methods that judge a label by its nearest neighbours' labels did worst under noise, because the distance is computed over general-purpose representations full of irrelevant features.
% - The adaptive hybrid representation selects task-relevant features from several representations; it is tested as a model input, beyond small molecules, and as the space for similarity and noise handling.
% - The aim throughout: property prediction that stays reliable when the training labels are noisy.
% ---- END TAKEAWAYS ----
% ---- NOTES: Introduction ----
% Takeaway: property prediction filters which compounds get tested, and its training labels are noisy in ways most studies do not model.
% Takeaway: the first half asks what survives label noise; the second half adapts the representation itself.
%
% Paragraph 2: label noise in molecular data
% - Add one sentence: this noise is rarely homoscedastic Gaussian. Its size varies with source and assay, as the pIC$_{50}$, pK$_i$ and single-assay figures in this paragraph show.
% - Change: "Measurements from a single assay lowers" becomes "lower".
%
% Paragraph 3: the first research question and the gap it fills
% - Change: "QSAR modelings" becomes "QSAR modelling".
% - Change: \cite{Cortes2015} becomes \citet{Cortes2015}, so the sentence reads with the author name.
%
% Paragraph 4: from noise robustness to adaptive representations
% - Change: noise mitigation is a subset of learning with noisy labels (LNL). "learning with noisy labels and noise mitigation" becomes "learning with noisy labels, including noise mitigation".
% - Say what was tried: 12 LNL methods adapted from classification to regression.
% - Several of them judge a label by the labels of its nearest neighbours, using a distance between molecules.
% - Change: "all of the methods that used similarity estimates failed" is too broad. The heteroscedastic GP compares molecules through its kernel and still held up.
% - Say instead: the methods that judged a label by its neighbours' labels did worst. The least noise-robust configurations all used a distance. Leave the counts for the next section.
% - Change: "vector operations" becomes "distance measures", naming the three tried: Tanimoto on ECFP4, Euclidean on PDV, and Mahalanobis.
% - Change: "representations" becomes the five tried: ECFP4, Sort \& Slice, PDV, and canonical and randomised SMILES.
% - For the similar property principle, cite \citet{Johnson1990}, as the Background does, or keep \citep{LopezPerez2024} as a review.
% - Change: typos. "simiarity" becomes "similarity"; "who's" becomes "whose"; "how it improve" becomes "how it improves".
% - Change: the last sentence says the hybrid improves performance. Chapters 5 and 6 test whether it does; their results are not yet in.
%
% Paragraph 5 (to write): the overall aim, then the hand-over. Three to four sentences.
% - The aim is property prediction that stays reliable when training labels are noisy.
% - The two parts: choosing models and representations that survive label noise, then adapting the representation itself.
% - Why it matters: property predictions filter which compounds are tested, so noise-driven errors waste experiments.
% - Hand-over line: the rest of the chapter states the questions, then outlines the thesis chapter by chapter.
% - Other theses in the field end their opener this way: one aim paragraph, then "The rest of this thesis is organised as follows".
% ---- END NOTES ----

New avenues in the drug discovery pipeline have opened up thanks to advances in large dataset collection, analysis, and computation. Therapeutic targets are identified, potential drug candidates are screened and optimized, and the most promising of these go on to pre-clinical and clinical trials \citep{Lill2007}. There are various methods by which potential drug candidates can be screened. Virtual screening allows for the rapid and large-scale analysis of libraries of compounds, and is a particularly useful approach used to speed up the drug discovery process \citep{Neves2018}. One standard screening method is the prediction of bioactivities and pharmacological properties of such candidates. There is a long history of using Quantitative Structure-Activity Relationship (QSAR) models to make these predictions, to act as a filter for compounds with undesirable properties, and thus reduce the time and cost associated with experimental screening \citep{Sabando2021}. QSAR models aim to predict the relationship between a molecule's molecular features as encoded within a particular molecular representation and a desired molecular property \citep{Sabando2021} of interest. A QSAR model typically consists of a molecular representation paired with a machine learning (ML) model. An ideal molecular representation will allow a model to learn only the relevant information about each molecule with respect to the target property. Representations are either `static' hand-crafted descriptors or circular fingerprints that encode predefined structural and physicochemical features \citep{rogers2010, David2020, sns, avalon}; or learned representations, where the encoding is derived from a predictive model, often taken from the activations of the final layer of a neural network, or a latent space embedding \citep{Sabando2021, kishimoto2023}. 

QSAR models are of course, sensitive to the quality of the input data on which they are trained. Experimental noise can arise from measurement errors, variability under experimental conditions, varying provenance of the data, and intrinsic biological and chemical variability. Even the very measures of model performance may be flawed if test and validation sets are laden with noise. \citet{Kolmar2021} show that the variance associated with experimental data can contribute more to prediction error than the model’s own error. There is also a limit to how much noise the model can tolerate; \citet{jorner2021} found that once the model error approaches the magnitude of the experimental measurement error on labels, adding more training data yields diminishing returns in predictive performance. In the context of activity prediction, the inherent variability within experimental setups and assays, batch effects, measurement errors, and biological variability introduce additional layers of complexity \citep{Kolmar2021, Fourches2010}. All of this increases the likelihood of encountering noise in the data \citep{Fang2022}. Even the molecular representations in the training data themselves may be incorrect or misleading \citep{Walters2023}, and chemical datasets contain too few repeat measurements of each compound to characterize the distribution the label is drawn from. Thus, models treat a single label as the population mean \citep{Kolmar2021}. Repeat measurements of the same compound-target pair disagree by about $0.68$ log units for pIC$_{50}$, and by about $0.54$ log units for pK$_i$ \citep{Kalliokoski2013, Kramer2012}. Measurements from a single assay lowers the mean disagreement from $0.50$ to $0.27$ log units \citep{landrum2024}.  

% TODO: rephrase: However, all of the methods that used simiarity metrics failed, even when I tried using different molecular representations and vector operations
This thesis addresses the challenge of conducting molecular property prediction when the underlying experimental data contain noisy labels. My initial research question is which combinations of molecular representations and predictive model architectures maintain predictive performance under label corruption, and whether probabilistic methods that estimate uncertainty justify their additional complexity. As I will show, both model and representation choice dominate predictive performance on clean data, however, this relationship inverts under label corruption as predictive architecture becomes the primary factor in determining noise robustness.  Most existing work on noise in QSAR modelings assumes that label noise can be modeled as random error drawn from a Gaussian distribution, and they evaluate either molecular representations or predictive model architectures in isolation. \cite{Cortes2015} examined 12 algorithms across multiple noise levels but used fixed representations. Studies that compare representations are typically done on clean data. No work has systematically compared the relative importance of representation versus architecture choice for noise robustness, assessed whether probabilistic methods justify their computational overhead in noisy settings, or tested whether findings generalize across different corruption patterns and molecular property types.

Building on these findings, I began exploring techniques involving learning with noisy labels and noise mitigation. Many of these methods incorporate measures of molecular similarity, relying on the assumption that molecules which are similar in structure share similar properties \citep{LopezPerez2024}. However, all of the methods that used simiarity estimates failed, even when I tried varying the molecular representations and vector operations from the standard ECFP4 and Tanimoto coefficient. I hypothesized that to generate accurate similarity estimates, you need a specialised molecular representation who's features relate directly to similarity in that molecular domain. This brings me to my next research question, how do you select the best set of features for a given prediction task? I will outline how I developed an adaptive feature selection methodology that takes in a set of molecular representations, identifies each representation's most important features, and combines a select set of those top features into a new hybrid representation. I will then demonstrate how it improve performance across a range of predictive tasks beyond the small molecule space, and how it can be applied to the context of similarity estimation and downstream noise mitigation methods. 
% TODO: how do I end this? Check other theses


\section{The questions this thesis asks}\label{sec:intro:questions}

% ---- NOTES: The questions this thesis asks ----
% Takeaway: four questions about QSAR models under label noise, all answered in Chapter 3.
% Takeaway: then questions about adapting the representation, each answered in a named chapter.
%
% Paragraph 1: the four questions of the first half
% - Add one closing sentence: Chapter~\ref{chap:robustness} answers all four.
% - Change: the last sentence introduces NoiseInject again. Chapter~\ref{chap:background} already describes it; point there.
%
% Paragraph 2: the questions of the second half, with the chapter that answers each
% - Add a question for Chapter~\ref{chap:flaws}: do standard representations, and the similarity computed over them, fail on a given dataset?
% - "How do you select the best set of features": Chapter~\ref{chap:method}, which gives the search and when and why the hybrid helps.
% - "Do hybrid representations achieve higher noise robustness": no chapter yet reports the retention comparison it asks for. Chapter~\ref{chap:similarity} reports Spearman at two noise levels, not the retention area.
% - Change: if that question stays, "(retention percentage)" becomes "normalised area under the retention curve", the robustness measure of Chapter~\ref{chap:robustness}.
% - "Do the selected hybrid features help stabilize distance-based noise mitigation": Chapter~\ref{chap:similarity}. Widen it to noise detection and estimation as well, since that chapter reports both.
% - "Does feature importance weighting improve": during hybrid construction and model training, Chapter~\ref{chap:method}. In the distance, and under noise, Chapter~\ref{chap:similarity}.
% - Change: in the importance-weighting question, "noise mitigation tasks" becomes "distance-based noise detection and learning with noisy labels".
% - "Can sample weighting or label smoothing provide superior noise handling": Chapter~\ref{chap:similarity}. Reword it.
% - Why reword: weighting and smoothing are reported only for cleaning real labels with nothing injected. No experiment applies them under Gaussian corruption.
% - Under Gaussian noise, Chapter~\ref{chap:similarity} instead surveys twelve LNL methods, on each downstream model's top two single representations, the hybrid and the importance-weighted hybrid.
% - Wording: "learning with noisy labels" for the broad class. Sample weighting and label smoothing are noise mitigation, a subset of it.
% - Add the similarity question: is predicting each molecule's activity from its nearest neighbours more accurate in a hybrid than in its base representations? (Chapter~\ref{chap:similarity}.)
% ---- END NOTES ----

% TODO: review, needs manual review
The first half of this thesis addresses four questions about the behavior of QSAR models under label noise. First, how does the choice of molecular representation and model architecture explain the variance in predictive performance, and does that division change when you instead consider how much performance a QSAR model is able to retain under increasing label noise? Second, what does the choice of model provide with label noise? Within this question we compare the effects of deterministic and probabilistic modeling. Third, does the type of label noise matter, or just the quantity of it? We explore different techniques of generating and adding label noise to both clean data, as well as data with inherent experimental noise. Fourth, what is the impact of artificial label noise on a model's uncertainty estimation? Does the estimated uncertainty rise with increased label noise and can these uncertainty values point to which labels are corrupted? Together, these questions aim to identify what, if anything, makes a QSAR model robust to noise. Finally, we introduce NoiseInject, an open-source Python package for benchmarking noise robustness and uncertainty quantification in machine learning models.

% TODO: review, needs manual review
The second half asks whether the representation itself can be adapted to the task:
\begin{enumerate}
    \item How do you select the best set of features for a given prediction task?
    \item Do hybrid representations achieve higher noise robustness (retention percentage) than their base representations?
    \item Do the selected hybrid features help stabilize distance-based noise mitigation?
    \item Does feature importance weighting improve effectiveness of hybrid representation creation, model training, and noise mitigation tasks?
    \item When sample removal is suboptimal, can sample weighting or label smoothing provide superior noise handling under Gaussian corruption?
\end{enumerate}

\section{From noise robustness to adaptive representations}

% ---- NOTES: From noise robustness to adaptive representations ----
% Takeaway: the bridge between the two halves, with the evidence behind each step.
% Takeaway: three short paragraphs; no method detail, each step points to the chapter that holds it.
%
% Paragraph 1: what the first half found (Chapter~\ref{chap:robustness})
% - Robustness is the normalised area under the R^2 retention curve as noise rises.
% - Under Gaussian noise on QM9, model architecture explains 49.5% of the variance in robustness. Representation explains 9.2%.
% - For R^2 at noise level 1.0 (noise one spread of the clean training labels), the shares are comparable: 29.1% for architecture, 26.6% for representation.
% - These R^2 shares are under noise, not on clean labels. Do not call them the clean-label split.
% - Result needed: the same variance split on R^2 before any noise is added. No claim about clean labels until it exists.
% - Systematic noise (shared within a scaffold family, or clipped at an assay limit) hurt models most. The shape of the noise distribution barely mattered.
% - Under Gaussian noise on QM9, mean predicted uncertainty rose with noise for every probabilistic model and representation.
% - Per-molecule uncertainty did not reliably pick out which labels were corrupted (Section~\ref{sec:uq}).
% - NGBoost and the Gaussian process fit a per-molecule scale. They gave the strongest per-molecule correlation between predicted uncertainty and injected noise.
% - Under censoring, models were most confident on the labels that had been clipped.
% - Consequence: uncertainty alone cannot say which labels to fix. That leads to methods that judge each label another way.
%
% Paragraph 2: the exploratory study of learning with noisy labels, and its negative result (Section~\ref{sec:lnl_motivation})
% - Learning with noisy labels (LNL) is the broad family. Noise mitigation, which removes, reweights or corrects labels, is one part of it.
% - Twelve LNL methods were adapted to regression and tested in 221 configurations on the QM9 HOMO-LUMO gap with Gaussian noise.
% - Several variants judged a label by its neighbours' labels, using a distance between molecules.
% - All 20 configurations with the lowest retention of R^2 under noise used such a distance.
% - This held for each distance tried: Tanimoto on ECFP4, Euclidean on physicochemical descriptors, and Mahalanobis.
% - The methods that held up estimated noise from the model's own loss or uncertainty, not from neighbours.
% - One dataset and one noise type: state it as the motivation, not a general result.
%
% Paragraph 3: why the distance failed, and the step to the hybrid (Chapters~\ref{chap:flaws}, \ref{chap:method} and \ref{chap:similarity})
% - Hypothesis: on a given dataset, most features of a general-purpose representation do not vary or do not relate to the property.
% - Every such feature adds a term to the distance, so irrelevant features dominate it. The growth is linear in the number of features, not exponential.
% - Chapter~\ref{chap:flaws} measures this on base representations and on distances.
% - The proposed fix: select a fixed number of task-relevant features from several representations, and compute distances in that compact space.
% - A first version showed that these selected feature sets also work as ordinary model inputs, with no distance involved.
% - That turned the question into how to choose the features, which is Chapter~\ref{chap:method}.
% - Chapter~\ref{chap:similarity} then returns to distances, noise detection and LNL, including noise mitigation, with the hybrid as the feature space.
% ---- END NOTES ----
\todo[inline, color=yellow!40]{Write this section from the notes above. The raw material is Section~\ref{sec:lnl_motivation}.}


\section{Outline of the thesis and contributions}

% ---- NOTES: Outline of the thesis and contributions ----
% Takeaway: one short paragraph per chapter: what it does, what it found (Chapters 3 and 4) or will test (Chapters 5 and 6), and where it is published or planned.
%
% Paragraph 1: Chapter 2 (Background)
% - Introduces, once, what the later chapters rely on: molecular representations, model families, molecular similarity and activity cliffs, noise in molecular data, and the evaluation protocol.
% - Defines the seven noise conditions used in Chapter 3.
% - Ends by introducing NoiseInject, the open-source benchmarking framework released with this work.
%
% Paragraph 2: Chapter 3 (Noise robustness across molecular representations and model architectures)
% - Adapted from the paper "NoiseInject: A Benchmarking Framework to Study Noise Robustness in QSAR Models" by Punt, Hanser, Werner and Morris, in preparation for the Journal of Cheminformatics. Say this in the paragraph.
% - Main task: the HOMO-LUMO gap on QM9. Its labels are computed, so they carry negligible measurement noise.
% - Repeated on three assay datasets with experimental labels: LogD, Caco-2 and hERG Ki.
% - Splits the variance in performance, and in how much performance is retained under noise, between representation and model architecture.
% - Compares seven noise conditions at a matched amount of noise.
% - Compares each model with its probabilistic counterpart, and asks what label noise does to predicted uncertainty.
% - Findings in words only, no numbers:
% - Model architecture matters more than representation for noise robustness. For R^2 under noise, the two matter about equally.
% - The shape of the noise distribution makes little difference. Noise shared within a scaffold family, and labels clipped at an assay limit (censoring), do the most damage.
% - NGBoost and random forests were the most robust models on QM9. Random forests were the most consistent across all four datasets.
% - Predicted uncertainty rises with noise on average, but does not pick out which labels were corrupted.
%
% Paragraph 3: Chapter 4 (The flaws of the similarity assumption)
% - A motivation chapter. It introduces no method, and its measurements are smaller than Chapter 3's.
% - First argument: on a given dataset, standard representations carry many features that do not vary, and some map different molecules to the same vector.
% - Second argument: similarity computed over those representations inherits both problems.
% - Third argument: methods that rely on that similarity fail as a result.
% - These include cleaning labels by comparing each molecule's label with its nearest neighbours' labels, a noise mitigation method tested in the exploratory study of learning with noisy labels. Say plainly that this study comes from the confirmation report.
% - Ends by stating what a representation would need to fix this. That is what Chapter 5 builds.
%
% Paragraph 4: Chapter 5 (Building adaptive hybrid representations)
% - Planned as the first of two method papers. Intended for submission; venue to be confirmed.
% - The method: from several base representations, it selects each one's most important features for the task.
% - It splits a fixed budget of features between the representations and combines the selected features into one representation.
% - Also describes the datasets and the experimental setup.
% - Will test when the hybrid helps, against the best single representation and against an ensemble of one model per representation.
% - Will test why it helps, against PCA, mRMR, recursive feature elimination and random features at the same number of features.
% - Will test whether the selected features can be read against the chemistry, and how stable the selection is across folds.
% - Will test whether search settings found on one dataset carry to another.
% - Will test whether it extends beyond small molecules: antibodies in most depth, and reactions, nucleotide sequences and cancer genomics as proofs of concept.
% - Will test whether it extends to data held at separate sites that cannot be pooled.
% - Result needed: one sentence of findings for Chapter 5.
%
% Paragraph 5: Chapter 6 (Hybrid representations for similarity and label noise)
% - Planned as the second method paper. Intended for submission; venue to be confirmed.
% - Returns to the failures in Chapter 4 and asks whether a hybrid representation repairs them.
% - Similarity: hybrids selected for predicting each molecule's activity from its five nearest neighbours, compared with standard representations.
% - Noise detection and estimation: whether neighbours in a hybrid pick out corrupted labels better than neighbours in a standard representation.
% - Learning with noisy labels, including noise mitigation, one noise type at a time.
% - Result needed: one sentence of findings for Chapter 6.
%
% Paragraph 6: Chapter 7 (Discussion and future directions)
% - Summarises each chapter.
% - Closing argument: model the noise before choosing a method. No single method fits every dataset, and simple methods are often as good.
% - Then limitations and future directions.
%
% Paragraph 7: Contributions (short: publications and software)
% - Publications: the Chapter 3 paper (in preparation) and the planned papers from Chapters 5 and 6. Point to the Research Papers list in the Summary of Research Contributions rather than repeating it.
% - Software: NoiseInject (open source), and the KIRBy software repository that implements the hybrid search and its experiments.
% - Repositories, licences and data: point to Appendix~\ref{app:availability}.
% ---- END NOTES ----
\todo[inline, color=yellow!40]{Write Sec 1.5 from the outline bullets above.}

\section{Collaborations and funding}

% ---- NOTES: Collaborations and funding ----
% Takeaway: who funded the work, what Lhasa Limited and OPIG contributed, and what is the author's own.
%
% Paragraph 1: funding
% - The DPhil was funded by the University of Oxford's UK EPSRC Centre for Doctoral Training in Sustainable Approaches to Biomedical Science: Responsible and Reproducible Research, grant EP/S024093/1.
% - It was also supported by Lhasa Limited, a not-for-profit organisation.
%
% Paragraph 2: Lhasa Limited
% - Thierry Hanser and Stephane Werner at Lhasa Limited co-supervised the work, with Garrett Morris as supervisor.
% - Chapter~\ref{chap:robustness} is adapted from the jointly authored NoiseInject manuscript (Punt, Hanser, Werner and Morris). Each author's contribution is in the Statement of Originality.
% - Name Lhasa's federated-learning platform, Effiris, in one clause \citep{Bassani2023}.
% - Lhasa supplied the hERG data (7,772 training and 4,383 test molecules) used in Chapter~\ref{chap:method} as data held at separate sites, the setting federated learning addresses \citep{Bassani2023, Hanser2023}.
%
% Paragraph 3: the Oxford Protein Informatics Group (OPIG)
% - OPIG pointed to antibody and protein datasets used in Chapters~\ref{chap:method} and~\ref{chap:similarity}. The full list is in Appendix~\ref{app:datasets}.
% - The collaboration also produced five posts on OPIG's blog, Blopig. Point to the Blogposts section rather than listing them.
%
% Paragraph 4: own work
% - One sentence: the adaptive hybrid representation, its software and the experiments in Chapters~\ref{chap:flaws} to~\ref{chap:similarity} are the author's own work.
% ---- END NOTES ----
\todo[inline, color=yellow!40]{Write Sec 1.6.}

\pagebreak
